\documentclass[a4paper, serif, 10pt]{moderncv}

%% moderncv
% style and color
\moderncvstyle{banking}
\moderncvcolor{black}

%% page layout/margins
\usepackage[top=2.5cm, bottom=2.5cm, left=1.5cm, right=1.5cm]{geometry}

\nopagenumbers{}

%% Babel config for Norwegian language
% ref:
% - https://latex3.github.io/babel/guides/locale-norwegian.html
\usepackage[norsk]{babel}

%% fontspec
\usepackage{fontspec}

\IfFontExistsTF{Linux Libertine}{
  \setmainfont[Ligatures={TeX, Common}, Numbers=OldStyle]{Linux Libertine}
}{}
\IfFontExistsTF{Linux Libertine O}{
  \setmainfont[Ligatures={TeX, Common}, Numbers=OldStyle]{Linux Libertine O}
}{}

%% CTeX
% CJK support, used to show CJK characters in the resume
%
% - fontset=none: disable builtin fontset but instead set the CJK font manually
% - heading=false: disable ctex heading
% - punct=kaiming: use kaiming punctuations styles for CJK
% - scheme=plain: use plain scheme, do not override \normalsize font size
% - space=auto: space settings for CJK characters
%
% ref:
% - http://ctan.mirrorcatalogs.com/language/chinese/ctex/ctex.pdf
\usepackage[UTF8, heading=false, punct=kaiming, scheme=plain, space=auto]{ctex}

\IfFontExistsTF{Noto Serif CJK SC}{
  \setCJKmainfont{Noto Serif CJK SC}
}{}
\IfFontExistsTF{Noto Sans CJK SC}{
  \setCJKsansfont{Noto Sans CJK SC}
}{}

%% line spacing
\usepackage{setspace}
\setstretch{1.125}

%% URL styling - use normal text instead of monospace
\urlstyle{same}

% Auto-underline all links
% Defer until \begin{document} so hyperref is already loaded
\AtBeginDocument{%
  \let\oldhref\href
  \renewcommand{\href}[2]{\oldhref{#1}{\underline{#2}}}%
}

%% Patch \httplink and \httpslink to handle full URLs
% The original moderncv macros always prepend http:// or https:// to URLs.
% This causes issues when the URL already has a protocol (e.g., https://example.com)
% resulting in malformed URLs like https://https://example.com.
% This patch checks if the URL already contains "://" and uses it directly if so.
% Note: We use \str_set:Nx (x-expansion) to expand macros like \@homepage first.
\ExplSyntaxOn
\str_new:N \g_yamlresume_url_str

\RenewDocumentCommand{\httpslink}{O{}m}{
  \str_set:Nx \g_yamlresume_url_str {#2}
  \str_if_in:NnTF \g_yamlresume_url_str {://}
    { \url{#2} }
    { \url{https://#2} }
}

\RenewDocumentCommand{\httplink}{O{}m}{
  \str_set:Nx \g_yamlresume_url_str {#2}
  \str_if_in:NnTF \g_yamlresume_url_str {://}
    { \url{#2} }
    { \url{http://#2} }
}
\ExplSyntaxOff

\name{Ingrid Sørensen}{}
\title{UX-designer med lidenskap for brukervennlige digitale løsninger}
\phone[mobile]{+47 912 34 567}
\email{ingrid@ingridsorensen.no}
\homepage{https://ingridsorensen.no}

\address{Storgata 23, Oslo, Oslo, Norge, 0180}{}{}

\extrainfo{{\small \faLinkedin}\ \href{https://www.linkedin.com/in/ingridsorensen}{@ingridsorensen} {} {} {} • {} {} {}
{\small \faBehance}\ \href{https://www.behance.net/ingridsorensen}{@ingridsorensen} {} {} {} • {} {} {}
{\small \faDribbble}\ \href{https://dribbble.com/ingridsorensen}{@ingridsorensen}}

\begin{document}

\maketitle

\section{Grunnleggende informasjon}

\cvline{}{\begin{itemize}
\item UX-designer med 6 års erfaring fra produktutvikling og digitale tjenester
\item Sterk kompetanse innen brukerresearch, prototyping, og brukertesting
\item Erfaring med å lede designprosesser fra idé til ferdig produkt
\item Tverrfaglig samarbeid med utviklere, produktledere og forretningssiden
\item Gode ferdigheter i Figma, Sketch og designsystemer
\end{itemize}}
\section{Utdanning}

\cventry{aug. 2014–juni 2017}
        {Bachelor, Interaksjonsdesign, Poeng: A / 5}
        {OsloMet – storbyuniversitetet}
        {\url{https://www.oslomet.no/}}
        {}
        {\begin{itemize}
\item Fordypet meg i brukersentrerte designprosesser og metodikk
\item Praksisprosjekt med oppdragsgiver fra offentlig sektor
\item Utvekslingssemester ved Københavns Universitet
\end{itemize}
\textbf{Kurs}: Brukersentrert design, Webutvikling og universell utforming, Informasjonsarkitektur, Designmetodikk, Kognitiv psykologi, Multimodal kommunikasjon}
\section{Arbeidserfaring}

\cventry{jan. 2021–Nå}
        {Senior UX-designer}
        {Nordlys Digital}
        {\url{https://www.nordlysdigital.no}}
        {}
        {\begin{itemize}
\item Leder designprosesser for kunder innen bank, helse og offentlig sektor
\item Utvikler designsystemer og gjenbrukbare UI-komponenter i Figma
\item Gjennomfører brukertester og dybdeintervjuer for å validere designbeslutninger
\item Samarbeider tett med utviklingsteam i Scrum-metodikk
\item Mentor for juniordesignere og holder interne workshops
\end{itemize}
\textbf{Nøkkelord}: Designsystem, Brukertesting, Figma, Scrum}

\cventry{feb. 2018–des. 2020}
        {UX/UI-designer}
        {Apphuset AS}
        {\url{https://www.apphuset.no}}
        {}
        {\begin{itemize}
\item Designet brukervennlige mobilapper og web-løsninger for startups og etablerte selskaper
\item Utførte brukerresearch, brukerreiser, wireframes og interaktive prototyper
\item Etablerte og vedlikeholdt designbibliotek i Figma og Sketch
\item Bidro i kundeoppfølging og presentasjon av designløsninger
\end{itemize}
\textbf{Nøkkelord}: Mobilapp, UI-design, Prototyping, Sketch}

\cventry{juli 2017–des. 2017}
        {UX-designpraktikant}
        {Kommunal IT}
        {\url{https://www.kommunal-it.no}}
        {}
        {\begin{itemize}
\item Assisterte senior designere i kartlegging av brukerbehov og utforming av innsynsløsninger
\item Produserte wireframes og klikkbare prototyper i Adobe XD
\item Deltok i arbeidsmøter med innbyggere og saksbehandlere
\end{itemize}
\textbf{Nøkkelord}: Adobe XD, Innsyn, Brukerbehov}
\section{Språk}

\cvline{Norsk}{Morsmål eller tospråklig nivå \hfill \textbf{Nøkkelord}: Morsmål}
\cvline{Engelsk}{Fullt profesjonelt nivå \hfill \textbf{Nøkkelord}: IELTS 8.0}
\cvline{Spansk}{Begrenset arbeidskunnskap \hfill \textbf{Nøkkelord}: Selvstudium}
\section{Ferdigheter}

\cvline{UX-design}{Ekspert \hfill \textbf{Nøkkelord}: Brukerresearch, Brukerreiser, Informasjonsarkitektur, Brukertesting, Interaksjonsdesign}
\cvline{UI-design}{Avansert \hfill \textbf{Nøkkelord}: Figma, Sketch, Adobe XD, Designsystem, Prototyping}
\cvline{Nettutvikling}{Begynner \hfill \textbf{Nøkkelord}: HTML, CSS, JavaScript}
\cvline{Prosjektledelse}{Viderekommen \hfill \textbf{Nøkkelord}: Scrum, Kanban, JIRA, Miro}
\section{Utmerkelser}

\cventry{nov. 2022}
        {Nasjonal UX-forening}
        {UX-prisen for inkluderende design}
        {}
        {}
        {\begin{itemize}
\item Mottatt for arbeidet med et inkluderende innloggingssystem for offentlige tjenester.
\end{itemize}}

\cventry{juni 2017}
        {OsloMet}
        {Årets studentprosjekt}
        {}
        {}
        {\begin{itemize}
\item Vant med bacheloroppgaven "Universell utforming av smarte hjem-løsninger".
\end{itemize}}
\section{Sertifikater}

\cventry{mars 2023}
        {Google}
        {Google UX Design Professional Certificate}
        {\url{https://coursera.org/verify/ux-design}}
        {}
        {}

\cventry{sep. 2021}
        {Nielsen Norman Group}
        {NN/g UX Certification}
        {\url{https://www.nngroup.com}}
        {}
        {}
\section{Publikasjoner}

\cventry{aug. 2022}
        {Brukerinvolvering i praksis: En håndbok for team}
        {Fagbokforlaget}
        {\url{https://www.fagbokforlaget.no}}
        {}
        {\begin{itemize}
\item Praktisk guide for hvordan team kan involvere brukere i hele produktutviklingsprosessen.
\item Medforfatter med eksempler fra offentlig sektor.
\end{itemize}}
\section{Prosjekter}

\cventry{jan. 2022–des. 2022}
        {Digital veiviser som hjelper innbyggere med å finne riktige offentlige tjenester.}
        {Innbyggerlos}
        {\url{https://www.innbyggerlos.no}}
        {}
        {\begin{itemize}
\item Initierte og ledet designarbeidet for en ny digital veiviser for offentlige tjenester.
\item Gjennomførte dybdeintervjuer og brukertester med 30+ innbyggere.
\item Endte opp med en løsning som reduserte henvendelser til kundesenteret med 25 \%.
\end{itemize}
\textbf{Nøkkelord}: Offentlig sektor, Tjenestedesign, Inkluderende design}

\cventry{feb. 2020–juni 2020}
        {Mobilapp for klimavennlige forbruksvalg.}
        {Klimasmart hverdag}
        {\url{https://www.klimasmart.no}}
        {}
        {\begin{itemize}
\item Designet konseptet for en app som gir forbrukere konkrete klimatips.
\item Bygde interaktive prototyper og utførte brukertester.
\item Vant intern innovasjonskonkurranse og ble videreutviklet av produktselskap.
\end{itemize}
\textbf{Nøkkelord}: Bærekraft, Mobilapp, Konseptutvikling}
\section{Interesser}

\cvline{Friluftsliv}{Fottur, Ski, Kajakk}
\cvline{Kreativitet}{Digital tegning, Keramikk, Fotografi}
\section{Frivillig arbeid}

\cventry{jan. 2020–juni 2023}
        {Designmentor}
        {Jentetech}
        {\url{https://www.jentetech.no}}
        {}
        {\begin{itemize}
\item Veileder unge jenter som vil lære seg UX-design.
\item Holder workshops i Figma og porteføljeutvikling.
\end{itemize}}


\cventry{aug. 2018–des. 2019}
        {Nettsidedesigner}
        {Frivilligsentralen}
        {\url{https://www.frivillig.no}}
        {}
        {\begin{itemize}
\item Redesignet nettsiden for en lokal frivilligsentral.
\item Sørget for universell utforming og forbedret innholdspresentasjon.
\end{itemize}}

\end{document}